\chapter{DESENVOLVIMENTO}
Neste capítulo deve ser descrito, de forma detalhada e estruturada, todo o processo de concepção, projeto e construção da solução proposta ou da investigação técnica realizada. O objetivo é permitir que o leitor compreenda as decisões de projeto e como os objetivos específicos foram executados na prática.

No contexto de Análise e Desenvolvimento de Sistemas, recomenda-se organizar este capítulo em seções que contemplem, conforme a natureza do projeto:
\begin{uemaball}
    \item \textbf{Levantamento e Análise de Requisitos:} especificação dos requisitos funcionais, requisitos não funcionais e regras de negócio do sistema;
    \item \textbf{Arquitetura e Tecnologias Utilizadas:} descrição da arquitetura adotada (ex.: MVC, microsserviços, cliente-servidor) e das linguagens, \textit{frameworks}, bibliotecas e banco de dados empregados;
    \item \textbf{Modelagem do Sistema:} apresentação e explicação dos diagramas (ex.: casos de uso, classes, entidade-relacional, fluxo de processos) que representam a estrutura e o comportamento da solução;
    \item \textbf{Prototipação e Implementação:} detalhamento das principais funcionalidades desenvolvidas, telas, integrações e algoritmos relevantes.
\end{uemaball}

Este é um exemplo de uso de notas de rodapé. O sistema foi desenvolvido utilizando a arquitetura de microsserviços\footnote{Para mais informações acesse: \url{https://aws.amazon.com/pt/microservices/}} para garantir maior escalabilidade.

Sempre que inserir ilustrações (diagramas, fluxogramas, capturas de tela ou gráficos), mencione-as previamente no texto (ex.: conforme ilustrado na Figura \ref{fig:bpmn-modelagem}), inserindo o título na parte superior e a fonte na parte inferior, de acordo com as normas da ABNT e da UEMA:

\begin{figure}[htb]
    \centering
    \caption{\label{fig:bpmn-modelagem}Representação da estrutura do Spring Data}
    \includegraphics[width=0.85\linewidth]{img/bpmn_modelagem.png}
    \fonte{\citeonline{ribeiro2015beneficios}}
\end{figure}
